\chapter{Results and Analysis}
\section{Aggregate Classification Performance}
Table~\ref{tab:overall} summarizes the complete balanced matrix. Mean accuracy
falls from \CleanAccuracy{} with original training labels to \NoisyAccuracy{}
with noisy labels, a loss of approximately 10.52 percentage points. Confidence
removal recovers \ConfidenceRemoveGain{} points, while combined removal recovers
\CombinedRemoveGain{} points. Neither restores the clean reference on average.
The corresponding macro-F1 and AUROC values also favor confidence removal
among the evaluated cleaning methods.

\reporttable{overall}{Overall test performance across the balanced matrix.}{tab:overall}
{Each method has 384 rows. Means weight datasets, models, noise types, rates,
and seeds equally. ``pp'' denotes percentage points. Repeated clean values
are a reference, not independent measurements.}

Relabeling is beneficial on average for confidence, combined, and disagreement
scoring, but their removal variants achieve higher aggregate accuracy.
Margin-only removal is the worst intervention: accuracy is 0.7939, below the
noisy baseline by 5.96 points. Margin relabeling also reduces accuracy, by
2.56 points. Thus applying a nominally plausible uncertainty score is not
sufficient to improve robustness; the relationship between the score and the
\emph{observed label} matters.

\section{RQ1: Sensitivity to Noise}
Table~\ref{tab:loss} separates the accuracy loss by dataset, noise mechanism,
and rate. The clean and noisy models are paired before aggregation. On Dry
Bean, the mean loss under symmetric 40\% corruption is 2.92 points, compared
with 20.75 points under cyclic corruption. Digits shows the same contrast:
8.82 versus 29.26 points. This suggests that a concentrated wrong-class
transition is more damaging than dispersing errors among many alternatives
for these multiclass tasks. It does not establish a universal ordering of
noise mechanisms, because only one cyclic map is evaluated.

\reporttable{noise_loss}{Accuracy loss from clean to noisy training, in percentage points.}{tab:loss}
{Each cell averages the paired loss across three models and four seeds.
Rates are nominal training-label corruption fractions.}

For Breast Cancer, both mechanisms are binary flips, and the differences
reflect sampled rows and per-class versus global corruption rather than
different destination labels. Wine is especially sensitive at high noise
rates: its small training and test partitions limit precision. One changed
test prediction on Wine corresponds to $1/36$, or about 2.78 accuracy points
in a single run. Small differences should not be interpreted as strong
rankings of methods.

Figures~\ref{fig:symmetric} and \ref{fig:cyclic} show selected conditions
across rates. They display averages, not uncertainty bands. The complete
release contains per-model plots with seed-standard-deviation shading;
those shaded regions are not confidence intervals over independent datasets.

\begin{figure}[htbp]
\centering\includegraphics[width=\linewidth]{accuracy_symmetric.pdf}
\caption{Accuracy under symmetric corruption. Each point averages three
models and four seeds within a dataset. Selected methods are shown for
legibility; all methods appear in Table~\ref{tab:overall}. Lines connect the
four evaluated rates and do not represent fitted continuous trends.}
\label{fig:symmetric}
\end{figure}
\begin{figure}[htbp]
\centering\includegraphics[width=\linewidth]{accuracy_class_dependent.pdf}
\caption{Accuracy under cyclic class-dependent corruption, aggregated as in
Fig.~\ref{fig:symmetric}. High-rate cyclic corruption is particularly damaging
on the multiclass tasks.}
\label{fig:cyclic}
\end{figure}

\subsection{Variation across Model Families}
Table~\ref{tab:models} reports dataset--model means across rates, mechanisms,
and seeds. RF has the strongest mean noisy-label accuracy on Wine and Digits;
LR has the strongest on Dry Bean. HGB has lower mean noisy accuracy than both
other models on each of the four datasets. This is evidence about the fixed
settings used here, not about optimally tuned boosting in general.
For Wine, HGB rises from 0.7439 to 0.8429 after confidence removal but remains
below the cleaned RF value of 0.9418.

\reporttable{models}{Mean test accuracy by dataset and classifier.}{tab:models}
{LR: logistic regression; RF: random forest; HGB: histogram gradient boosting.
Each cell averages two noise types, four rates, and four seeds. ``Conf.'' and
``Comb.'' denote confidence and combined scoring; rem./rel. denote removal/relabeling.}

Removal is not always preferable to relabeling. On Digits with RF, confidence
relabeling reaches 0.9370 versus 0.9330 for removal. On Dry Bean with HGB,
the corresponding means are nearly tied. A global average therefore does
not justify selecting one action blindly for every model--dataset pair.

\section{RQ2: Detection of Injected Errors}
Table~\ref{tab:detection} counts each detector once per experimental combination,
not once for both removal and relabeling. Confidence has the highest average
detection F1, 0.8072, followed by combined, 0.7504, disagreement, 0.6827, and
margin, 0.2425. Because the selected fraction closely matches the injected
fraction, precision and recall are nearly equal; this is partly a design
property rather than independent evidence from two unrelated criteria.

\reporttable{detection}{Injected-label-error detection, averaged over 384 combinations per detector.}{tab:detection}
{Detection is evaluated on training rows. Removal and relabeling share each
detector's mask and are not double-counted.}

\begin{figure}[htbp]
\centering\includegraphics[width=\linewidth]{detection.pdf}
\caption{Detection F1 by nominal rate and corruption mechanism. Each point
averages the four datasets, three models, and four seeds. No test labels are
used to compute a suspiciousness score.}
\label{fig:detection}
\end{figure}

Margin-only selection behaves differently from confidence selection.
A confidently wrong observed label can receive a low margin score when the
model strongly predicts another class. Conversely, a correctly labeled but
ambiguous example can receive a high margin score. The poor detection and
downstream outcomes are consistent with this mismatch between uncertainty
and label error. This is a mechanism-based interpretation of the scores,
not a causal experiment isolating boundary examples.

\section{RQ3: Effects of Cleaning and Relabeling}
Figure~\ref{fig:gains} shows the paired accuracy gain within each dataset.
Confidence removal improves all four dataset-level averages, from 2.24
points on Dry Bean to 7.87 points on Wine. Margin removal reduces all four
averages, from 4.14 points on Dry Bean to 7.12 points on Digits. This pattern
is not merely the result of pooling the largest dataset with smaller ones:
datasets receive equal weight and their separate effects are visible.

\begin{figure}[htbp]
\centering\includegraphics[width=\linewidth]{dataset_gains.pdf}
\caption{Paired accuracy gains over noisy training, in percentage points.
Every cell averages 24 settings and four seeds for one dataset. Numeric
annotations preserve interpretation without color.}
\label{fig:gains}
\end{figure}

Table~\ref{tab:inference} reports both effect size and the limits of inference.
Confidence removal improves 94.8\% of the 96 seed-averaged settings. This is
not the same as its proportion of individually improved seed runs: aggregation
reduces some seed variability before a setting is counted as a win.
The confidence-removal exploratory interval is entirely positive, yet the
exact two-sided $p$-value is 0.125 and its Holm-adjusted value is 1.000.
There is no contradiction: the bootstrap percentile interval and discrete
exact test answer the uncertainty question differently with very sparse
dataset-level evidence. The interval must not be used to claim significance.

\reporttable{inference}{Paired accuracy effects and exploratory dataset-level inference.}{tab:inference}
{Gain and interval endpoints are percentage points. Intervals resample four
dataset means. Win rates use 96 seed-averaged settings. $p_H$ is Holm-adjusted
within the eight accuracy comparisons. None is significant at 0.05.}

\subsection{Does the Combined Score Help?}
The direct comparison in Table~\ref{tab:ablation} gives combined removal an
average deficit of 0.58 accuracy points relative to confidence removal;
combined relabeling has a deficit of 0.68 points. Combined scoring exceeds
disagreement and margin on the grand mean, but does not beat the best simple
component. All reported adjusted comparisons remain nonsignificant.

\reporttable{ablation}{Combined scoring minus each component detector under the same action.}{tab:ablation}
{Positive gain favors combined. Each comparison uses matched
dataset--model--noise--rate--seed observations; inference uses four dataset
means. Holm adjustment covers the six accuracy comparisons.}

Breast Cancer's combined and confidence downstream results coincide in the
aggregate, as the binary ranking analysis predicts. For each of the three
multiclass datasets, confidence has a larger mean gain for both actions.
The combination can shift highly uncertain examples upward without providing
additional evidence that their observed labels are wrong. A different weight
choice might behave differently, but choosing it after inspecting these
test results would require a new, independently evaluated selection procedure.
The present negative result should not be obscured by retrospective tuning.

\subsection{A Concrete Failure despite Good Detection}
The saved-model example uses Dry Bean, RF, symmetric 20\% noise, confidence
scoring, and seed 42. Table~\ref{tab:diagnostic} makes the failure case explicit.
Detection F1 is 0.9151, but removal and relabeling both reduce test accuracy
relative to the noisy baseline. They improve macro AUROC in this setting,
showing that probability ranking and hard-label accuracy can move in
different directions. The small differences do not support a separate
statistical conclusion from one seed.

\reporttable{diagnostic}{Illustrative saved-model setting: Dry Bean/RF/symmetric 20\%/seed 42.}{tab:diagnostic}
{Confidence is used for both cleaning actions. PR AUC is the trapezoidal area
under the micro-averaged precision--recall curve, not average precision.}

The noisy classifier's confusion matrix in Fig.~\ref{fig:confusion} identifies
where residual errors occur. Class IDs are defined in Chapter~4. This example
is a diagnostic illustration, not an unbiased choice of the best-performing
setting. The full matrix, rather than this single case, supports the aggregate
comparison. ROC and PR curves for the same saved classifier are provided in
Appendix~A.

\begin{figure}[htbp]
\centering\includegraphics[width=.88\linewidth]{noisy_baseline_confusion_matrix.pdf}
\caption{Test confusion matrix for the illustrative noisy RF classifier.
Rows are reference classes and columns predicted classes. The test set has
2,709 examples and was not corrupted.}
\label{fig:confusion}
\end{figure}

\section{Threats to Validity}
\textbf{Internal validity.} Shared splits, training-only out-of-fold scoring,
pipeline-local scaling, exact-count corruption, and duplicate removal reduce
leakage risk. Saved checksums and re-execution provide evidence that the
tables correspond to the code. Nevertheless, testing cannot exclude every
implementation error. The same integer seed drives multiple random processes;
results use deterministic paired conditions rather than independently
randomized substreams. A future design could study sensitivity to separate
split, corruption, and estimator seeds.

\textbf{Construct validity.} The injected mask is an operational definition
of noise. It does not establish that the original labels are correct. The
cyclic transition is arbitrary in class-ID space, and the known-rate budget
is unavailable in many applications. Margin measures uncertainty, while
confidence combines possible label error with model misspecification.
The reserved validation data are unused; this avoids tuning leakage but
does not evaluate practical threshold selection.

\textbf{External validity.} Four numeric datasets, only one binary task,
and no real audited annotation errors limit generalization. Digits is
image-derived and the other datasets are relatively small. There are no
mixed categorical tables, deployment shifts, feature-noise experiments,
or calibrated estimates of label-error prevalence. No clinical or other
high-stakes use is supported by these experiments.
Random stratified partitions do not establish generalization to unseen
patients, writers, or collection sites; subject-group independence is not
evaluated by this benchmark.

\textbf{Comparative validity.} Only fixed estimator parameters and four
simple scores are compared. Confident learning, robust losses, stronger
tabular learners, instance-dependent corruption, and adaptive budgets remain
unmeasured. No runtime benchmark establishes that the proposed pipeline is
faster than alternatives. The combined score's modest engineering complexity
is not evidence of a new algorithmic contribution.

\textbf{Statistical validity.} The 3,840 rows share datasets, splits,
corruptions, and repeated references. They cannot be treated as independent
replicates. Four dataset means are insufficient for precise population-level
inference; the convenience sample also weakens generalization beyond the
observed tasks. Retrospective best-by-setting tables are descriptive only
and are not an evaluated model-selection policy.
